\section{Cómo funcionan la entrada y la salida estándar}
En todos los problemas del curso, las soluciones se entregan como un programa que lee desde la entrada estándar (\textit{standard input}) y escribe en la salida estándar (\textit{standard output}). No se reciben argumentos por línea de comandos ni se abren archivos desde el código: el programa siempre lee y escribe usando los mecanismos estándar del lenguaje elegido (por ejemplo, \texttt{Scanner}/\texttt{System.in} en Java, \texttt{input()} en Python, \texttt{cin} en C++, o \texttt{readline} en JavaScript).

\subsection{Redirección}
Para probar un programa sin escribir cada caso a mano, la consola permite redirigir un archivo hacia la entrada estándar y capturar la salida estándar en otro archivo, sin modificar una sola línea del programa:

\begin{itemize}
    \item \texttt{<} redirige el contenido de un archivo hacia la entrada estándar del programa.
    \item \texttt{>} redirige la salida estándar del programa hacia un archivo.
\end{itemize}

\subsection{Comando de ejecución según el lenguaje}
Suponiendo una solución en un archivo llamado \texttt{Problema}, con los casos de prueba en \texttt{entrada.in} y la salida esperada en \texttt{salida.out}:

\begin{verbatim}
java Problema < entrada.in > salida.out
python Problema.py < entrada.in > salida.out
./Problema < entrada.in > salida.out
node Problema.js < entrada.in > salida.out
\end{verbatim}

Java y C++ deben compilarse antes de ejecutarse; cómo hacerlo depende del entorno de desarrollo que cada estudiante elija usar.
